\chapter{Monitoring et Logging}

\section{Presentation}

Le monitoring repose sur trois piliers complementaires~:
\begin{itemize}
    \item \textbf{CloudWatch} --- Monitoring AWS natif (logs, metriques, alarmes)
    \item \textbf{Prometheus + Grafana} --- Metriques applicatives et dashboards personnalises
    \item \textbf{SonarQube} --- Qualite du code et analyse statique
\end{itemize}

\section{CloudWatch}

\subsection{Logs}

Les conteneurs ECS emettent leurs logs via le driver \texttt{awslogs}
vers le groupe \texttt{/ecs/deus-dashboard-backend} avec une retention de 14 jours.
L'agent CloudWatch (installe par Ansible) collecte egalement les logs de l'agent ECS.

\subsection{Alarmes}

\begin{table}[H]
    \centering
    \begin{tabularx}{\textwidth}{lXl}
        \toprule
        \textbf{Alarme} & \textbf{Condition} & \textbf{Action} \\
        \midrule
        ECS CPU elevee & CPU > 80\% pendant 10 min & Email SNS \\
        ECS memoire elevee & Memoire > 80\% pendant 10 min & Email SNS \\
        ALB erreurs 5xx & > 10 erreurs en 5 min & Email SNS \\
        RDS CPU elevee & CPU > 80\% pendant 10 min & Email SNS \\
        RDS stockage bas & < 2 Go restants & Email SNS \\
        \bottomrule
    \end{tabularx}
    \caption{Alarmes CloudWatch configurees}
\end{table}

\section{Prometheus}

Prometheus est deploye sur une instance EC2 dediee au monitoring
(\texttt{34.207.159.92}), accessible a l'adresse
\texttt{http://34.207.159.92:9090}.
Il collecte les metriques via le protocole pull (scraping)~:

\begin{table}[H]
    \centering
    \begin{tabularx}{\textwidth}{llX}
        \toprule
        \textbf{Job} & \textbf{Cible} & \textbf{Metriques} \\
        \midrule
        deus-backend & ALB:80/api/metrics & Requetes HTTP (duree, compteur), cache hits/miss, heap Node.js \\
        node-exporter & localhost:9100 & CPU, memoire, disque, reseau de l'instance \\
        prometheus & localhost:9090 & Metriques internes de Prometheus \\
        \bottomrule
    \end{tabularx}
    \caption{Jobs de scraping Prometheus}
\end{table}

\subsection{Metriques applicatives}

Le backend expose un endpoint \texttt{/api/metrics} au format Prometheus
grace a la librairie \texttt{prom-client}~:

\begin{lstlisting}[language=bash,caption={Metriques exposees par le backend}]
# Duree des requetes HTTP (histogramme)
http_request_duration_seconds_bucket{method,route,status_code}

# Compteur total de requetes
http_requests_total{method,route,status_code}

# Hits et miss du cache DynamoDB
cache_hit_total{route}
cache_miss_total{route}

# Metriques Node.js par defaut
nodejs_heap_size_used_bytes
nodejs_active_handles_total
process_cpu_user_seconds_total
\end{lstlisting}

\section{Grafana}

Grafana est deploye sur la meme instance que Prometheus
(\texttt{http://34.207.159.92:3000}) et fournit des dashboards interactifs.

\subsection{Sources de donnees}

\begin{itemize}
    \item \textbf{Prometheus} --- Metriques applicatives et systeme
    \item \textbf{CloudWatch} --- Metriques AWS (ECS, ALB, RDS)
\end{itemize}

\subsection{Dashboard Backend Overview}

Le dashboard pre-provisionne comprend~:
\begin{itemize}
    \item Taux de requetes par seconde (par methode, route, code de statut)
    \item Latence des requetes (percentiles p50 et p95)
    \item Taux de hit du cache DynamoDB (jauge)
    \item Taux d'erreurs 5xx
    \item Utilisation CPU et memoire de l'instance (via Node Exporter)
    \item Utilisation du heap Node.js
\end{itemize}

\subsection{Acces au monitoring}

\begin{table}[H]
    \centering
    \begin{tabular}{lll}
        \toprule
        \textbf{Service} & \textbf{URL} & \textbf{Port} \\
        \midrule
        Grafana & \texttt{http://34.207.159.92:3000} & 3000 \\
        Prometheus & \texttt{http://34.207.159.92:9090} & 9090 \\
        Node Exporter & \texttt{http://34.207.159.92:9100} & 9100 \\
        \bottomrule
    \end{tabular}
    \caption{URLs d'acces aux services de monitoring}
\end{table}

% TODO: Capture d'ecran du dashboard Grafana
% \begin{figure}[H]
%     \centering
%     \includegraphics[width=\textwidth]{grafana-dashboard.png}
%     \caption{Dashboard Grafana --- Backend Overview}
% \end{figure}

\section{SonarQube --- Qualite du code}

SonarQube analyse le code source pour detecter~:
\begin{itemize}
    \item Les bugs potentiels
    \item Les vulnerabilites de securite
    \item Le code duplique
    \item Les code smells (mauvaises pratiques)
    \item La couverture de tests
\end{itemize}

\subsection{Integration CI}

L'analyse SonarQube s'execute automatiquement dans le pipeline CI
(GitHub Actions) apres les etapes de build. Le fichier
\texttt{sonar-project.properties} configure les sources a analyser~:

\begin{lstlisting}[language=bash,caption={Configuration SonarQube}]
sonar.projectKey=deus-dashboard
sonar.sources=src,backend/src
sonar.exclusions=**/node_modules/**,**/dist/**
\end{lstlisting}

\subsection{Developpement local}

SonarQube est egalement disponible en local via Docker Compose
sur le port 9000 (\texttt{http://localhost:9000}).

% TODO: Capture d'ecran SonarQube
% \begin{figure}[H]
%     \centering
%     \includegraphics[width=\textwidth]{sonarqube-dashboard.png}
%     \caption{Dashboard SonarQube}
% \end{figure}

\section{Architecture du monitoring}

\begin{lstlisting}[language={},caption={Flux de monitoring}]
   Backend (ECS) --scrape--> Prometheus
                 --awslogs--> CloudWatch
                 --sonar-scan--> SonarQube (CI)

   Node Exporter --scrape--> Prometheus

   Prometheus --> Grafana (dashboards)
   CloudWatch --> Grafana (datasource)
   CloudWatch --> SNS --> Email (alarmes)
\end{lstlisting}
